% Chunk: physical PDF pages 471--474 (printed pages 448--451).
% Semantic coverage: conclusion of Section 10.4; Section 10.5 through the
% sentence continuing onto physical page 475.
% Equations: (10.4.28), (10.5.1)--(10.5.18).
% Figures: none. Footnotes: one. Uncertainties: none.

For a renormalized Dirac field, Eqs.~\eqref{eq:10.3.31} and
\eqref{eq:10.4.27} tell us that on the mass shell
\begin{equation}
  \bar u_{k'}\Gamma^\mu(k,k)u_k=\bar u_{k'}\gamma^\mu u_k,
  \tag{10.4.28}\label{eq:10.4.28}
\end{equation}
where $[i\gamma_\mu k^\mu+m]u_k=[i\gamma_\mu k'^\mu+m]u_{k'}=0$.
Thus the renormalization of the fermion field ensures that the radiative
corrections to the vertex function $\Gamma_\mu$ cancel when a fermion on the
mass shell interacts with an electromagnetic field with zero momentum
transfer, as is the case when we set out to measure the fermion's electric
charge. If we had not used a renormalized fermion field then the corrections
to the vertex function would have just cancelled the corrections due to
radiative corrections to the external fermion lines, leaving the electric
charge again unchanged.

\subsection{Gauge Invariance}

The conservation of electric charge may be used to prove a useful result for
the quantities
\begin{align}
  M^{\mu\mu'\cdots}_{\beta\alpha}(q,q',\ldots)
  \equiv{}&\int d^4x\int d^4x'\cdots
  e^{-iq\cdot x}e^{-iq'\cdot x'}\cdots \notag\\
  &\quad\cdot
  \OutBra{\beta}
  T\{J^\mu(x),J^{\mu'}(x'),\ldots\}
  \InKet{\alpha}.
  \tag{10.5.1}\label{eq:10.5.1}
\end{align}
In theories like spinor electrodynamics in which the electromagnetic
interaction is linear in the field $A^\mu$, this is the matrix element for
emission (and/or absorption) of on- or off-shell photons having four-momenta
$q$, $q'$, etc. (and/or $-q$, $-q'$, etc.), with external line photon
coefficient functions or propagators omitted, in an arbitrary transition
$\alpha\to\beta$. Our result is that Eq.~\eqref{eq:10.5.1} vanishes when
contracted with any one of the photon four-momenta:
\begin{equation}
  q_\mu M^{\mu\mu'\cdots}_{\beta\alpha}(q,q',\ldots)
  =q'_{\mu'}M^{\mu\mu'\cdots}_{\beta\alpha}(q,q',\ldots)
  =\cdots=0.
  \tag{10.5.2}\label{eq:10.5.2}
\end{equation}
Since $M$ is defined symmetrically with respect to the photon lines, it will
be sufficient to show the vanishing of the first of these quantities.

For this purpose, note that by an integration by parts
\begin{align}
  q_\mu M^{\mu\mu'\cdots}_{\beta\alpha}(q,q',\ldots)
  ={}&-i\int d^4x\int d^4x'\cdots
  e^{-iq\cdot x}e^{-iq'\cdot x'}\cdots \notag\\
  &\quad\cdot
  \OutBra{\beta}
  \frac{\partial}{\partial x^\mu}
  T\{J^\mu(x),J^{\mu'}(x'),\ldots\}
  \InKet{\alpha}.
  \tag{10.5.3}\label{eq:10.5.3}
\end{align}
The electric current $J^\mu(x)$ is conserved, but this does not immediately
imply that Eq.~\eqref{eq:10.5.3} vanishes, because we still have to take
account of the $x^0$-dependence contained in the theta functions that appear
in the definition of the time-ordered product. For instance, for just two
currents
\begin{equation*}
  T\{J^\mu(x)J^\nu(y)\}
  =\theta(x^0-y^0)J^\mu(x)J^\nu(y)
  +\theta(y^0-x^0)J^\nu(y)J^\mu(x)
\end{equation*}
so, taking account of the conservation of $J^\mu(x)$:
\begin{align}
  \frac{\partial}{\partial x^\mu}T\{J^\mu(x)J^\nu(y)\}
  ={}&\delta(x^0-y^0)J^0(x)J^\nu(y)
  -\delta(y^0-x^0)J^\nu(y)J^0(x) \notag\\
  ={}&\delta(x^0-y^0)[J^0(x),J^\nu(y)].
  \tag{10.5.4}\label{eq:10.5.4}
\end{align}
With more than two currents, we get an equal-time commutator like this
(inside the time-ordered product) for each current aside from $J^\mu(x)$
itself. To evaluate this commutator, we recall that (as shown in the previous
section) for any product $F$ of field operators and their adjoints and/or
derivatives
\begin{equation*}
  [J^0(\mathbf{x},t),F(\mathbf{y},t)]
  =-q_F F(\mathbf{x},t)\delta^3(\mathbf{x}-\mathbf{y}),
\end{equation*}
where $q_F$ is the sum of the $q_\ell$s for the fields and field derivatives
in $F$, minus the sum of the $q_\ell$s for the field adjoints and their
derivatives. For the electric current, $q_J$ is zero; $J^\nu(y)$ is itself an
electrically neutral operator. It follows that
\begin{equation}
  [J^0(\mathbf{x},t),J^\nu(\mathbf{y},t)]=0
  \tag{10.5.5}\label{eq:10.5.5}
\end{equation}
and therefore Eq.~\eqref{eq:10.5.4} vanishes, so that
Eq.~\eqref{eq:10.5.3} gives
\begin{equation}
  q_\mu M^{\mu\mu'\cdots}_{\beta\alpha}(q,q',\ldots)=0
  \tag{10.5.6}\label{eq:10.5.6}
\end{equation}
as was to be proved.

There is an important qualification here. In deriving Eq.~\eqref{eq:10.5.5}
we should take into account the fact that a product of fields at the same
spacetime point $y$ like the current operator $J^\nu(y)$ can only be properly
defined through some regularization procedure that deals with the infinities
in such products. In many cases it turns out that there are non-vanishing
contributions to the commutator of $J^0(\mathbf{x},t)$ with the regulated
current $J^\nu(\mathbf{y},t)$, known as Schwinger terms~\hyperref[ref:10.6]{[6]}.
Where the current includes terms arising from a charged scalar field $\Phi$,
there are additional regulator-independent Schwinger terms involving
$\Phi^\dagger\Phi$. However, all these Schwinger terms are cancelled in
multi-photon amplitudes by the contribution of additional interactions that
are quadratic in the electromagnetic field, either arising from the regulator
procedure (if gauge-invariant) or, as for charged scalars, directly from terms
in the Lagrangian. We will be dealing mostly with charged spinor fields, and
will use a regularization procedure (dimensional regularization) that does not
lead to Schwinger terms, so in what follows we will ignore this issue and
continue to use the naive commutation relation~\eqref{eq:10.5.5}.

The same argument yields a result like Eq.~\eqref{eq:10.5.2} even if other
particles besides photons are off the mass shell, provided that all charged
particles are taken on the mass shell, i.e., kept in the states
$\OutKet{\beta}$ and $\InKet{\alpha}$. Otherwise, the
left-hand side of Eq.~\eqref{eq:10.5.2} receives contributions from
non-vanishing equal-time commutators, such as those we encountered in the
derivation of the Ward identity in the previous section.

One consequence of Eq.~\eqref{eq:10.5.2} is that $S$-matrix elements are
unaffected if we change any photon propagator $\Delta_{\mu\nu}(q)$ by
\begin{equation}
  \Delta_{\mu\nu}(q)\longrightarrow
  \Delta_{\mu\nu}(q)+\alpha_\mu q_\nu+q_\mu\beta_\nu
  \tag{10.5.7}\label{eq:10.5.7}
\end{equation}
or if we change any photon polarization vector by
\begin{equation}
  e_\rho(\mathbf{k},\lambda)\longrightarrow
  e_\rho(\mathbf{k},\lambda)+ck_\rho,
  \tag{10.5.8}\label{eq:10.5.8}
\end{equation}
where $k^0\equiv|\mathbf{k}|$, and $\alpha_\mu$, $\beta_\nu$, and $c$ are
entirely arbitrary (not necessarily constants, and not necessarily the same
for all propagators or polarization vectors.) This is (somewhat loosely)
called the gauge invariance of the $S$-matrix.

To prove this result it is only necessary to display the explicit dependence
of the $S$-matrix on photon polarization vectors and propagators
\begin{align}
  S_{\beta\alpha}\propto{}&
  \int d^4q_1\,d^4q_2\cdots
  \Delta_{\mu_1\nu_1}(q_1)\Delta_{\mu_2\nu_2}(q_2)\cdots \notag\\
  &\quad\cdot
  e^*_{\rho_1}(\mathbf{k}'_1,\lambda'_1)
  e^*_{\rho_2}(\mathbf{k}'_2,\lambda'_2)\cdots
  e_{\sigma_1}(\mathbf{k}_1,\lambda_1)
  e_{\sigma_2}(\mathbf{k}_2,\lambda_2)\cdots \notag\\
  &\quad\cdot
  M_{ba}^{\mu_1\mu_2\cdots\nu_1\nu_2\cdots
  \rho_1\rho_2\cdots\sigma_1\sigma_2\cdots}
  (-q_1,-q_2,\ldots,q_1,q_2,\ldots,
  -k'_1,-k'_2,\ldots,k_1,k_2,\ldots).
  \tag{10.5.9}\label{eq:10.5.9}
\end{align}
where $M^{\mu\sigma\cdots}$ is the matrix element~\eqref{eq:10.5.1}
calculated in the absence of electromagnetic interactions.\footnote{The
states $a$ and $b$ are the same as $\alpha$ and $\beta$, but with photons
deleted. Note that the arguments of $M$ are all taken to be \emph{incoming}
four-momenta, which is why we have to insert various signs for some of the
arguments of $M$ in Eq.~\eqref{eq:10.5.9}.} The invariance of
Eq.~\eqref{eq:10.5.9} under the `gauge transformations'~\eqref{eq:10.5.7}
and~\eqref{eq:10.5.8} follows immediately from the conservation
conditions~\eqref{eq:10.5.2}. (In Section 9.6 we used the path-integral
formalism to prove a special case of this theorem, that vacuum expectation
values of time-ordered products of gauge-invariant operators are independent
of the constant $\alpha$ in the propagator~\eqref{eq:9.6.21}.) This result is
not as elementary as it looks, as it applies not to individual diagrams, but
only to sums of diagrams in which the current vertices are inserted in all
possible places in the diagrams.

There is a particularly important application of Eq.~\eqref{eq:10.5.2} to the
calculation of the photon propagator. The complete photon propagator,
conventionally called $\Delta'_{\mu\nu}(q)$, takes the form
\begin{equation}
  \Delta'_{\mu\nu}(q)=\Delta_{\mu\nu}(q)
  +\Delta_{\mu\rho}(q)M^{\rho\sigma}(q)\Delta_{\sigma\nu}(q),
  \tag{10.5.10}\label{eq:10.5.10}
\end{equation}
where $M^{\rho\sigma}$ is proportional to the matrix element~\eqref{eq:10.5.1}
with two currents and $\alpha$ and $\beta$ both the vacuum state, and
$\Delta_{\mu\nu}$ is the bare photon propagator, written here in a general
Lorentz-invariant gauge as
\begin{equation}
  \Delta_{\mu\nu}(q)
  =\frac{\eta_{\mu\nu}-\xi(q^2)q_\mu q_\nu/q^2}{q^2-i\epsilon}.
  \tag{10.5.11}\label{eq:10.5.11}
\end{equation}
From Eq.~\eqref{eq:10.5.2} we have here $q^\mu M_{\mu\nu}(q)=0$, so that
\begin{equation}
  q^\mu\Delta'_{\mu\nu}(q)=q^\mu\Delta_{\mu\nu}(q)
  =\frac{q_\nu[1-\xi(q^2)]}{q^2-i\epsilon}.
  \tag{10.5.12}\label{eq:10.5.12}
\end{equation}
On the other hand, just as we did for scalar and spinor fields in Section
10.3, we may express the complete photon propagator in terms of a sum
$\Pi^*(q)$ of graphs with two external photon lines that (unlike $M$) are
one-photon-irreducible:
\begin{align}
  \Delta'(q)={}&\Delta(q)+\Delta(q)\Pi^*(q)\Delta(q)
  +\Delta(q)\Pi^*(q)\Delta(q)\Pi^*(q)\Delta(q)+\cdots \notag\\
  ={}&[\Delta(q)^{-1}-\Pi^*(q)]^{-1}.
  \tag{10.5.13}\label{eq:10.5.13}
\end{align}
or in other words
\begin{equation}
  \Delta'_{\mu\nu}(q)=\Delta_{\mu\nu}(q)
  +\Delta_{\mu\rho}(q)\Pi^{*\rho\sigma}(q)\Delta'_{\sigma\nu}(q).
  \tag{10.5.14}\label{eq:10.5.14}
\end{equation}
Then in order to satisfy Eq.~\eqref{eq:10.5.12}, we must have
\begin{equation}
  q_\rho\Pi^{*\rho\sigma}(q)=0.
  \tag{10.5.15}\label{eq:10.5.15}
\end{equation}
This together with Lorentz invariance tells us that $\Pi^*(q)$ must take the
form
\begin{equation}
  \Pi^{*\rho\sigma}(q)=(q^2\eta^{\rho\sigma}-q^\rho q^\sigma)\pi(q^2).
  \tag{10.5.16}\label{eq:10.5.16}
\end{equation}
Then Eq.~\eqref{eq:10.5.13} yields a complete propagator of the form
\begin{equation}
  \Delta'_{\mu\nu}(q)
  =\frac{\eta_{\mu\nu}-\bar\xi(q^2)q_\mu q_\nu/q^2}
  {[q^2-i\epsilon][1-\pi(q^2)]},
  \tag{10.5.17}\label{eq:10.5.17}
\end{equation}
where
\begin{equation}
  \bar\xi(q^2)=\xi(q^2)[1-\pi(q^2)]+\pi(q^2).
  \tag{10.5.18}\label{eq:10.5.18}
\end{equation}

Now, because $\Pi^*_{\mu\nu}(q)$ receives contributions only from
one-photon-irreducible graphs, it is expected not to have any pole at $q^2=0$.
(There is an important exception in the case of broken gauge symmetry, to be
discussed in Volume II.) In particular, the absence of poles at $q^2=0$ in the
$q_\mu q_\nu$ term in $\Pi^*_{\mu\nu}(q)$ tells us that the function
$\pi(q^2)$ in Eq.~\eqref{eq:10.5.16} also has no such pole, and so the pole in
the complete propagator~\eqref{eq:10.5.17} is
